%% Repère cartésien orthonormé (O ; x, y, z) en perspective cavalière.
%% z vertical, y vers la droite, x en fuite vers le bas-gauche.
%% Les trois vecteurs unitaires i, j, k sont en rouge (solideA).
\documentclass[tikz,border=4pt]{standalone}
\usepackage{liaisons}
\begin{document}
\begin{tikzpicture}[x=1cm,y=1cm, scale=1.12,
  axe/.style={-{Stealth[length=5pt]}, line width=0.6pt},
  unitaire/.style={-{Stealth[length=5.5pt]}, line width=1.6pt, solideA}]

% --- les trois axes -------------------------------------------------
\coordinate (O) at (0,0);
\draw[axe] (O) -- (0,2.5)    node[above=1pt, font=\small] {$z$};
\draw[axe] (O) -- (2.9,0)    node[right=1pt, font=\small] {$y$};
\draw[axe] (O) -- (215:2.3)  node[below left=-1pt, font=\small] {$x$};

% --- les vecteurs unitaires ----------------------------------------
\draw[unitaire] (O) -- (0,0.95)
  node[midway, left=2pt, font=\small, text=solideA] {$\vec{k}$};
\draw[unitaire] (O) -- (0.95,0)
  node[midway, above=2pt, font=\small, text=solideA] {$\vec{\jmath}$};
\draw[unitaire] (O) -- (215:0.95)
  node[midway, above left=-2pt, font=\small, text=solideA] {$\vec{\imath}$};

% --- origine --------------------------------------------------------
\fill (O) circle (1.5pt);
\node[font=\small] at (0.3,-0.32) {$O$};
\end{tikzpicture}
\end{document}
